% ============================================================================
\chapter{ELF 加载器}
\label{ch:elf-loader}
% Stage 11 — 待实现
% ============================================================================

\section{本章目标}

\begin{itemize}
  \item 理解 ELF32 文件格式（ELF header、Program header、Section header）
  \item 实现 ELF 加载器：解析 ELF header → 加载 \texttt{.text}/\texttt{.data} 到用户地址空间
  \item 用 VMA 记录各段映射（代码段 r-x、数据段 rw-）
\end{itemize}

\section{背景知识：ELF 格式}

% TODO:
% - ELF header: e_ident (magic 0x7F ELF), e_type (ET_EXEC), e_entry, e_phoff
% - Program Header: p_type (PT_LOAD), p_vaddr, p_filesz, p_memsz, p_flags
% - 只需处理 PT_LOAD 段

\subsection{ELF 加载流程}

% TODO:
% 1. 验证 ELF magic + e_type=ET_EXEC + e_machine=EM_386
% 2. 遍历 Program Headers
% 3. 对每个 PT_LOAD: 分配物理页 → 映射到 p_vaddr → 复制文件内容 → 注册 VMA
% 4. 设置入口点 = e_entry

\section{数据结构}

% TODO: Elf32_Ehdr, Elf32_Phdr 结构体定义

\section{实现详解}

% TODO: elf_load() 函数逐步讲解

\section{VMA 集成}

% TODO: 每个 PT_LOAD 段注册为一个 VMA，权限从 p_flags 映射

\section{验证}

% TODO: 加载一个最小 ELF (用交叉编译器编译) → 跳转执行 → syscall exit

\section{本章小结}

ELF 加载器让内核能运行外部编译的程序。下一章实现系统调用，
让用户态程序能请求内核服务。
